\section{Conclusiones}

La incorporación de emociones dentro de agentes inteligentes constituye una de las principales líneas de investigación en computación afectiva y sistemas multiagente. A través de distintos modelos psicológicos y arquitecturas cognitivas, numerosos trabajos han demostrado que las emociones pueden mejorar significativamente la adaptación, interacción y comportamiento social de los agentes artificiales.

En este trabajo se han revisado los principales enfoques utilizados para representar emociones computacionalmente, incluyendo modelos discretos, dimensionales e híbridos. Asimismo, se han analizado arquitecturas afectivas como EBDI y WASABI, así como propuestas relacionadas con emociones sociales y sistemas multiagente normativos.

A pesar de los avances alcanzados, continúan existiendo importantes desafíos relacionados con la subjetividad emocional, la adaptación cultural y la integración entre emoción, cognición y aprendizaje.

El creciente desarrollo de agentes conversacionales avanzados y sistemas de Inteligencia Artificial social hace prever que la computación afectiva seguirá adquiriendo una gran relevancia en los próximos años, especialmente en aplicaciones relacionadas con interacción humano-computador, robótica social y simulación social inteligente.